\documentclass[11pt]{beamer}
\usetheme{Singapore}
%\usecolortheme{dove}
\usepackage{etoolbox}
\usepackage{ragged2e}
\usepackage[utf8]{inputenc}
\usepackage[italian]{babel}
\usepackage{scrextend}
\usepackage{tikz}
\usepackage{adjustbox}
\usepackage{multirow}
\usepackage{media9}%

\AtBeginSection[]{\frame{\Large \centering \insertsection} \subsection{} }
%\AtBeginSection[]{\subsection{}}

\newcommand{\nologo}{\setbeamertemplate{logo}{}}

\apptocmd{\frame}{}{\justifying}{}

%\author{Cap. Alberto Azzalini}
\title[Analisi Forense]{Furti e truffe: come difenderci?}
%\subtitle{Introduzione alle modalità operative}

\logo{\includegraphics[width=1cm]{pics/logo.png}}

\institute{\includegraphics[width=2cm]{pics/logo.png}\\A cura della Compagnia Carabinieri di Vipiteno}

\date{8 aprile 2016}

%subject{}

\setbeamercovered{transparent}

%\setbeamertemplate{navigation symbols}{}

\begin{document}
	{
	\nologo
	\maketitle
	}
	\begin{frame}
		\frametitle{Indice dei contenuti}
		\tableofcontents
	\end{frame}

	\section{Furti in appartamento}
	\begin{frame}
		\frametitle{Le principali tipologie di furti in appartamento}
		\begin{itemize}
			\item per distrazione (porte o finestre lasciate aperte);
			\item con scasso (porte o finestre divelte con piede di porco, finestre infrante);
			\item con il buco (trapano a mano e attrezzo sagomato)
		\end{itemize}
	\end{frame}

	\begin{frame}
		\frametitle{I furti per distrazione}
		Sono i casi in cui i ladri non devono fare molta fatica per entrare in casa:
		\begin{itemize}
			\item finestre o porte-finestre aperte (anche ai piani alti);
			\item finestre aperte ``con anta a ribalta'';
			\begin{figure}
				\centering
				\includegraphics[width=0.6\linewidth]{pics/aprifinestre}
			\end{figure}

			\item porte di ingresso non chiuse a chiave;
			\item autovetture lasciate aperte con a bordo le chiavi di casa.
		\end{itemize}
	\end{frame}
	
	\begin{frame}
		\frametitle{Furti con scasso}
		Una finestra chiusa non è un grosso ostacolo per un ladro determinato:
		\begin{itemize}
			\item le finestre tradizionali si aprono in 20 secondi con un semplice cacciavite;
			\item se il ladro non teme di fare rumore (se la nostra casa è isolata o in una zona molto rumorosa), può anche decidere di spaccare semplicemente il vetro.
		\end{itemize}
		
		In questi casi solo le inferriate (o almeno le tapparelle abbassate) possono rappresentare un ostacolo al tentativo di furto.
	\end{frame}
	
	\begin{frame}
		\frametitle{Furti ``con il buco'' [1/2]}
		\begin{center}%
						\includemedia[
						width=8cm,
						totalheight=8cm,
						keepaspectratio,
						activate=pageopen,
						addresource=pics/scasso.mp4,
						playbutton= fancy,
						flashvars={source=pics/scasso.mp4
							& scaleMode=noScale}
						]{}{VPlayer.swf}
		\end{center}
	\end{frame}
		
	\begin{frame}
		\frametitle{Furti ``con il buco'' [2/2]}
		\begin{itemize}
			\item non serve un trapano elettrico;
			\begin{figure}
				\centering
				\includegraphics[width=0.7\linewidth]{"pics/trapano a mano"}
			\end{figure}
			\item non fa rumore;
			\item è veloce;
			\item occorre una finestra in legno o PVC senza anima in alluminio.

		\end{itemize}		
	\end{frame}
	
	\begin{frame}
		\frametitle{Furti}
		\framesubtitle{Come difendersi}
		\textbf{Non aiutiamo i ladri:} chiudiamo sempre finestre, imposte e tapparelle anche se usciamo solo per poco tempo! 

		\centering{\textbf{Un ladro esperto non ha bisogno di più di 10 minuti per razziare un appartamento.}}		
	\end{frame}

	\section[Truffe]{Truffe: agli anziani, ma non solo\dots}
	\begin{frame}
		\frametitle{Come avviene la truffa}
		Le cosiddette truffe agli anziani, che però colpiscono anche persone che anziane non sono, possono avvenire:
		\begin{itemize}
			\item presso l'abitazione;
			\item in strada (e poi eventualmente continuare in casa).
		\end{itemize}
		In entrambi i casi, i malviventi si basano sulla capacità di convincere le vittime della loro buona fede, facendo leva, genere, su:
		\begin{itemize}
			\item affetti;
			\item senso di colpa;
			\item dubbio di essere in pericolo.
		\end{itemize} 
	\end{frame}	
	\begin{frame}
		\frametitle{I falsi funzionari}
		\framesubtitle{Come riconoscerli}
		
		Forse è la più classica delle truffe: i malviventi si presentano presso l'abitazione di persone anziane con diversi pretesti, spacciandosi per:
				
		\begin{itemize}
			\item operatori di polizia;
			\item operatori di aziende di servizio (acqua, luce, gas, \dots);
			\item assistenti sociali;
			\item funzionari del catasto.
		\end{itemize}
		L'obiettivo, in questo caso, è prevalentemente il furto di denaro o beni preziosi.
	\end{frame}
	
	\begin{frame}
		\frametitle{I falsi funzionari}
		\framesubtitle{Come operano}
		Il truffatore a volte si presenta da solo, ma più spesso sono in due e agiscono in vari modi. Vediamo due esempi:
		\begin{itemize}
			\item entrano entrambi nell'appartamento e, mentre uno distrae l'anziano, l'altro cerca di rubare soldi o valori;
			\item entra in casa una persona singola con un pretesto e cerca di portare via un oggetto qualsiasi. 
			
			In un secondo momento il complice, \textbf{spacciandosi per un agente di polizia}, si presenta presso l'abitazione chiedendo alla persona anziana se riconosce l'oggetto e invita a controllare se manca altro. In questo modo, è la stessa persona che mostra al malintenzionato il posto dove conserva soldi o preziosi.
			
		\end{itemize}
	\end{frame}
	
	\begin{frame}
		\frametitle{I falsi funzionari}
		\framesubtitle{Cosa sapere e cosa fare}
		
		Qualche consiglio:
		
		\begin{itemize}
			\item prima di fare entrare in casa sconosciuti, accertarne l'identità: non sussiste alcun obbligo di farli entrare;
			\item in caso di dubbio, chiedere aiuto \textbf{prima} di fare entrare il sedicente funzionario. Quasi sempre è sufficiente dire che intendete chiamare i Carabinieri per verificare l'identità;
			\item fare le verifiche \textbf{prima} che la persona sia entrata in casa, in quanto non sappiamo come possano reagire;
			\item qualora i truffatori si trovino già in casa, è importante osservarli bene e rilevare eventuali particolari che possano aiutare a identificarli;
			\item chiamare le forze dell'ordine, \textbf{senza toccare nulla}, per permettere di rilevare eventuali impronte.
		\end{itemize}
		

		
	\end{frame}
	
	\begin{frame}
		\frametitle{Il prestito non restituito}
		\framesubtitle{Come agiscono i truffatori}
		Le vittime vengono avvicinate da persone estranee che:
		\begin{itemize}
			\item affermano di aver fatto un prestito ad un parente (magari un nipote) della vittima;
			\item pretendono di ottenere la restituzione del denaro come asseritamente promesso.
		\end{itemize}
		
		Facendo leva sul legame familiare la vittima viene indotta a pagare, a volte anche andando in banca o in posta per fare il necessario prelievo della somma richiesta che a volte è anche considerevole.
	\end{frame}
	
	\begin{frame}
		\frametitle{Il prestito non restituito}
		\framesubtitle{Come difendersi}
		Qualche consiglio:
		\begin{itemize}
			\item chiedere il nome del parente, \textbf{ma senza suggerirlo!}
			
			Evitare cose del tipo: ``Mio nipote? Ma quale, Mario?'';
			\item contattare subito il proprio parente per verificare la situazione o le forze dell'ordine;
			\item cercare di attirare l'attenzione di altre persone, alzando al voce e chiedendo spiegazioni;
			\item se si hanno dubbi, paura o non si riesce a chiamare nessuno, è meglio assecondare il malvivente e, giunti in banca o in posta -- che sono ambienti protetti --, farsi aiutare dall'operatore dello sportello per avvertire i parenti e/o le Forze dell'Ordine.
		\end{itemize}
	\end{frame}
	
	\begin{frame}
		\frametitle{Le truffe ``all'americana'' e lo ``scam nigeriano''}
		\framesubtitle{La falsa beneficienza o l'eredità inaspettata}
		La vittima viene contattata -- spesso per email -- da una persona che deve fare una donazione o provvedere a consegnare una eredità. 
		\begin{itemize}
			\item il malintenzionato però non possiede al momento l'anticipo per il notaio chiede alla vittima di anticipare la somma a fronte di una percentuale del ricavo finale dell'operazione;
			\item la vittima viene accompagnata a ritirare l'anticipo -- spesso si tratta di una somma consistente -- e poi, saliti sulla macchina dei truffatori, ci si reca dal notaio;
			\item durante il tragitto i malviventi si ricordano che per l'atto sicuramente servirà una marca da bollo, fermano presso un tabacchino e chiedono alla vittima di andarla ad acquistare;
			\item non appena il truffato scende dall'auto, i trufatori si danno alla fuga con il denaro appena prelevato.
		\end{itemize}

		
	\end{frame}
	\begin{frame}
		\frametitle{Le truffe ``all'americana'' e lo ``scam nigeriano''}
		\framesubtitle{Le false pietre preziose}
		Gli autori sono due: 
		\begin{itemize}
			\item una persona si finge (o è) uno straniero che per un'urgenza deve raggiungere il paese d'origine, ma non ha disponibilità di soldi liquidi per il viaggio. 	Costui ferma la sua vittima per strada e cerca di vendere un oggetto molto prezioso (8-10 mila euro) che, vista l'urgenza, è disposto a lasciare per un prezzo molto vantaggioso (1-2 mila euro);
			\item  un altro signore distinto che arriva lì ``per caso'' e che dice di essere un gioielliere: visti i preziosi, si offre di comprarle per 5mila euro;
		\end{itemize}
	
		
		
		Il presunto straniero però insiste che sia l'altra persona a comprarle, facendogli credere che sia un affare che per \textit{correttezza} non vuole concedere al gioielliere.
		
		
	\end{frame}
	\begin{frame}
		\frametitle{Le truffe ``all'americana'' e lo ``scam nigeriano''}
		\framesubtitle{Come difendersi}
		
		Il truffatore fa leva sulla convinzione della vittima di fare un affare e la presenza del gioielliere al quale, nonostante l'offerta maggiore, non vuole cedere i preziosi in quanto li ha proposti prima al malcapitato, rinforza tale convinzione.
		
		Ricordiamoci che:
		\begin{itemize}
			\item \textbf{nessuno regala nulla}, tantomeno uno sconosciuto per la strada;
			\item occorre riflettere bene sulle proposte che si ricevono in circostanze strane;
			\item è bene chiedere di dimostrare in qualche modo la veridicità della proposta, prendere tempo e consultarsi in ogni caso con una persona di propria fiducia.
		\end{itemize}	
	\end{frame}
	
	\begin{frame}
		\frametitle{Ricontrollare il denaro}
		
		Il truffatore segue la vittima fino a casa, \textbf{dopo che questi ha fatto un prelievo al bancomat}. Si presenta alla porta e:
		\begin{itemize}
			\item sostiene di essere un funzionario della banca o della posta;
			\item afferma che sia necessario verificare che, tra le banconote prelevate, non ce ne sia una di falsa;
			\item mentre controlla, le sostituisce effettivamente con banconote false.
		\end{itemize}
		
		\textbf{Nessun funzionario di banca o delle Poste verrà mai a casa per controllare le banconote!}
		
		Non si deve permettere alla persona di entrare in casa, avvisando immediatamente le Forze dell'Ordine.
		
	\end{frame}
	
	\begin{frame}
		\frametitle{Il vecchio amico o la giacca macchiata}
		Si tratta di due trucchi per effettuare un borseggio, più che vere truffe.
		\begin{itemize}
			\item i truffatori (che sono spesso donne con bambini), dopo aver urtato la vittima e avergli fatto cadere sulla giacca gelato o caffè, con la scusa di ripulirgliela, le rubano il portafogli;
			\item una persona ben vestita e dall'aria distinta si avvicina alla vittima, facendo finta di conoscerla e abbracciandola calorosamente. Il vero scopo è di sottrarre il portafogli.
		\end{itemize}
		\textbf{In questi casi non bisogna distrarsi e fare sempre mente locale su dove si trovano i soldi.}
	\end{frame}
%	\begin{frame}
%		\begin{center}%
%			\includemedia[
%			width=10cm,
%			activate=pageopen,
%			addresource=pics/flv.flv,
%			playbutton= fancy,
%			flashvars={source=flv.flv}
%			]{Movie}{VPlayer.swf}
%		\end{center}%
%	\end{frame}
\end{document}